% Standalone combined formula guide: compile this file as the main document.
% Sources: contents/introduction.tex, chapter2_shortened.tex,
%          chapter3_shortened.tex and chapter4_shortened.tex.
% SINGLE COMPLETE FILE. Supersedes and merges all of:
%   formula_guide_chapters1_2.tex, formula_guide_chapters1_4.tex,
%   formula_guide_chapter5.tex, formula_guide_chapter6.tex,
%   appendix_formulas.tex
% Upload this file alone and set it as the Overleaf main document.
% Nothing else is needed: the bibliography is embedded and there are no images.
% Bibliography is embedded below, so no separate .bib upload is required.
%
% NO-REPETITION RULE: each distinct relation is printed exactly once, at the
% earliest chapter that prints it in the thesis. Later chapters that restate it
% cross-reference the equation number instead of reprinting the algebra.
\documentclass[11pt,a4paper]{article}
\usepackage[T1]{fontenc}
\usepackage[utf8]{inputenc}
\usepackage[margin=25mm]{geometry}
\usepackage{amsmath,amssymb}
\usepackage[hidelinks]{hyperref}
\numberwithin{equation}{section}
\setlength{\parindent}{0pt}
\setlength{\parskip}{0.5em}
\allowdisplaybreaks
\title{Formula Guide: Complete\\\large All formulas in the thesis, Chapters 1 to 6}
\author{}
\date{}
\begin{document}
\maketitle

\textbf{Scope.} This guide covers \textbf{every formula in the thesis}. It
follows all six chapters in their reading order, using the shortened version of
each chapter where one exists. It covers every displayed equation and the inline
formulas, dimensions and ranges used in them. Repeated symbols are explained
with their formula rather than listed as separate formulas.

\textbf{No formula is printed twice.} Chapters 2, 3 and 4 restate several of
the same relations, by design: Chapter 2 develops a mechanism, Chapter 3
reports the literature's version of it and any divergence, and Chapter 4 fixes
the implemented configuration. Each distinct relation therefore appears here
exactly once, at the earliest chapter that prints it, and every later
restatement cross-references that equation number instead of reprinting the
algebra. Entries marked \textbf{No new formula} exist to record what the later
chapter adds --- a scope, a constraint, a divergence --- not to repeat it.

\textbf{Where the mathematics actually is.} It is very unevenly spread, and
that is worth knowing before reading. Chapter 2 carries 16 of the thesis's 25
displayed equation blocks, Chapter 3 carries 7 and Chapter 4 carries 2.
\textbf{Chapters 1, 5 and 6 contain no displayed mathematics at all} --- verified
by scanning each source for displayed-math delimiters and equation
environments. Their sections here exist to record that, to promote the
formula-like quantities they state only in prose, and to index the notation
they use.

\textbf{How to read the sources.} A citation identifies a supplied paper that
actually presents the equation or supports the stated mechanism. It does not
automatically identify the equation's original inventor. Standard definitions,
algebraic counts and project-specific adaptations are marked explicitly where
no originating paper has been verified in the supplied collection. Paper page
numbers refer to the supplied PDF, whose preprint pagination may differ from
the final publication. Equations added to express prose mathematically are
identified as explanatory expansions.

\textbf{Notation.} Symbols are local to each topic: $v$ is vertical displacement
in optical flow but variance in SimAM; $\alpha$ is magnification gain in EVM but
a class weight in focal loss; $\epsilon$ may mean a numerical stabiliser or a
label-smoothing strength. These quantities are not interchangeable.

\tableofcontents
\clearpage

\section{Chapter 1: Introduction}
\label{fg:chapter1}

\textbf{Source file:} \texttt{contents/introduction.tex}. No shortened
introduction exists in this checkout. The chapter contains \textbf{no explicit
inline or displayed mathematical formulas}. It describes the four binary
component switches, twelve valid configurations, 25 subject folds and the
156-clip working set in prose. It also introduces matched comparisons without
writing their calculation as an equation. There are therefore no Chapter 1
equations or equation-specific paper attributions to extract. Formal equations
from later chapters have not been inserted into this section.

\section{Chapter 2: Background}
\label{fg:chapter2}

\textbf{Source file:} \texttt{contents/chapter2\_shortened.tex}. Its phenomenon
section (2.1) contains definitions but no mathematical formulas. The formulas
begin with recording and input preparation. Chapter 2 contains sixteen displayed
equation blocks; the guide also explains its inline calculations and notation.

\subsection{Recording and input preparation: thesis section 2.2}

\subsubsection{Duration and the number of frame intervals}
\label{fg:capture}
\textbf{Thesis location:} 2.2.1, High-speed capture.

The chapter's inline expression is $0.3f_s$ for an event lasting $0.3$ seconds.
Its general form, written out here for explanation, is
\begin{equation}
    n_{\mathrm{intervals}}\approx \tau f_s.
    \label{fg:eq:capture}
\end{equation}
Here $\tau$ is duration in seconds and $f_s$ is frames per second. Multiplying
them gives a number of temporal intervals: $0.3\times200=60$, versus 7.5--9 at
25--30 fps. Exact recorded frame counts depend on endpoint alignment; intervals
and frames are not the same count. Subsampling also changes the spacing between
the images that reach the model.

\textbf{Source.} This is the definition of sampling rate, not a new equation
from a research paper. Yan et al.~\cite{yan2014}, apparatus description on PDF
pages 2--3, support the CASME II recording rate of 200 fps.

\subsubsection{Image dimensions}
\label{fg:geometry}
\textbf{Thesis locations:} 2.2.2--2.2.3.

The expressions $640\times480$ and $224\times224$ denote image width and height
in pixels. They are dimensions, not learned transformations. The recorded study
uses the original full frames and resizes them to a square; this changes aspect
ratio and does not perform face alignment or cropping.

\textbf{Source.} Yan et al.~\cite{yan2014}, PDF page 2, report the original
capture dimensions. The square resize and choice of raw frames are project
settings, verified against the working manifest and frame loader; they are not
implied by the database paper.

\subsection{Motion magnification: thesis section 2.3}

\subsubsection{Small-motion approximation}
\label{fg:evm}
\textbf{Thesis location:} 2.3.1, The Eulerian idea.

% SOURCE-DISPLAY-01
\begin{equation}
    I(x,t)=f(x+\delta(t))
       \approx f(x)+\delta(t)\frac{\partial f}{\partial x}.
    \label{fg:eq:evm-taylor}
\end{equation}
$I(x,t)$ is observed intensity at image position $x$ and time $t$; $f$ is the
reference image profile; $\delta(t)$ is a small displacement. The derivative
$\partial f/\partial x$ measures how rapidly intensity changes across position.
The equation says that moving a textured profile slightly changes its intensity
approximately in proportion to displacement. It drops higher-order terms, so
large displacements need not satisfy the approximation.

\textbf{Source.} Bai et al.~\cite{bai2021}, PDF page 2, section 3.1,
equation (2), present this first-order Taylor expansion. They explain an existing
Eulerian magnification formulation; this citation does not credit them with
inventing Taylor expansion or EVM.

\subsubsection{The amplified intensity signal}
% SOURCE-DISPLAY-02
\begin{equation}
    \widetilde I(x,t)\approx
       f(x)+(1+\alpha)\delta(t)\frac{\partial f}{\partial x}.
    \label{fg:eq:evm-gain}
\end{equation}
$\widetilde I$ is the reconstructed intensity and $\alpha$ is the gain applied
to the retained temporal variation. The factor $1+\alpha$ includes the original
motion term plus its amplified addition. Within the small-motion approximation,
this resembles displacement increased by that factor. It does not guarantee
that all amplified variation is useful facial motion.

\textbf{Source.} Bai et al.~\cite{bai2021}, PDF page 2, equations (3)--(4).
The thesis retains this rationale but implements a hard Fourier mask and shared
gain across Laplacian detail bands, rather than reproducing every filtering
choice described in that paper.

\subsubsection{Pyramid blur kernel and output range}
\label{fg:pyramid}
\textbf{Thesis locations:} 2.3.2 and 2.3.5.

The inline kernel is
\begin{equation}
    h=[1,4,6,4,1]/16.
    \label{fg:eq:blur}
\end{equation}
It is a five-tap weighted average: the centre receives the greatest weight and
the weights sum to one. Applying it separately along image rows and columns
smooths spatial detail before downsampling. The formula above is the blur
kernel; it is not a statement that every reconstruction operation uses the
same scaling.

\label{fg:clip}
After magnification the implementation clips intensities to $[0,255]$ and
converts them to 8-bit greyscale. This range is a storage constraint, not the
range of a physical displacement or a trained probability.

\textbf{Source.} Project implementation:
\texttt{Stage1\_DataPipeline/evm\_magnifier.py}. Bai et al.~\cite{bai2021}
support the Laplacian-pyramid mechanism, but the exact kernel and output
conversion are identified here as code choices, not attributed to that paper.

\subsection{Optical flow and strain: thesis section 2.4}

\subsubsection{Brightness constancy}
\label{fg:flow}
\textbf{Thesis location:} 2.4.2.
% SOURCE-DISPLAY-03
\begin{equation}
    I(x,y,t)=I(x+u\Delta t,y+v\Delta t,t+\Delta t).
    \label{fg:eq:brightness}
\end{equation}
The same moving image point is assumed to retain its intensity between two
times. Here $u$ and $v$ are horizontal and vertical velocities, so multiplication
by the time interval $\Delta t$ gives displacement. In the discrete project
tensors, the same symbols instead describe pixels of displacement between
selected frames; they must not be interpreted as calibrated velocity per second.

\textbf{Source.} OFF-ApexNet~\cite{liong2019a}, supplied PDF page 5,
equation (4) and the definitions immediately below it. That paper's notation is
equivalent after substituting its displacement variables. Its particular
onset--apex estimator is not the complete sequence-processing method of this thesis.

\subsubsection{The optical-flow constraint}
% SOURCE-DISPLAY-04
\begin{equation}
    I_xu+I_yv+I_t=0.
    \label{fg:eq:flow-constraint}
\end{equation}
$I_x$ and $I_y$ are spatial intensity derivatives and $I_t$ is the temporal
derivative. Taylor expansion of brightness constancy gives this constraint.
It is one equation with two unknown motion components; an edge constrains
motion along its intensity gradient but not along the edge itself. A flow
estimator therefore needs additional neighbourhood assumptions. Brightness
changes and large motion can violate the approximation.

\textbf{Source.} OFF-ApexNet~\cite{liong2019a}, PDF page 5,
equations (5)--(7). The thesis's use of Farneb\"ack flow is a separate estimator
choice, also reported for micro-expression recognition by Zhao
et al.~\cite{zhaoS2021}; the constraint alone does not specify that estimator.

\subsubsection{The small-strain tensor}
\label{fg:strain}
\textbf{Thesis location:} 2.4.3.

Write the displacement vector as $\mathbf{u}=[u,v]^{\mathsf T}$. Its gradient
contains the spatial derivatives of both components; superscript $\mathsf T$
means transpose. The symmetric part is
% SOURCE-DISPLAY-05
\begin{equation}
\begin{aligned}
    \varepsilon&=\tfrac12\left(\nabla\mathbf{u}
                         +(\nabla\mathbf{u})^{\mathsf T}\right),\\
    \varepsilon_{xx}&=\frac{\partial u}{\partial x},\qquad
    \varepsilon_{yy}=\frac{\partial v}{\partial y}.
\end{aligned}
    \label{fg:eq:strain-normal}
\end{equation}
% SOURCE-DISPLAY-06
\begin{equation}
    \varepsilon_{xy}=\varepsilon_{yx}
       =\tfrac12\left(\frac{\partial u}{\partial y}
                    +\frac{\partial v}{\partial x}\right).
    \label{fg:eq:strain-shear}
\end{equation}
$\varepsilon_{xx}$ and $\varepsilon_{yy}$ describe local extension or
compression; the equal off-diagonal terms describe shear. A uniform translation
has constant displacement and therefore zero strain. This does not confer
immunity to arbitrary head rotation, perspective effects or flow errors.

\textbf{Source.} Shreve et al.~\cite{shreve2011}, PDF page 2,
section III-B, equations (2)--(3). Although the paper uses the term finite
strain, it explicitly assumes small motion. The symmetric-gradient expression
shown here is the linearised, small-strain approximation, not a general
finite-deformation tensor.

\subsubsection{Implemented strain magnitude and the Frobenius comparison}
\label{fg:magnitude}
% SOURCE-DISPLAY-07
\begin{equation}
    \varepsilon_{\mathrm{mag}}
       =\sqrt{\varepsilon_{xx}^2+\varepsilon_{yy}^2
                    +\varepsilon_{xy}^2}.
    \label{fg:eq:strain-magnitude}
\end{equation}
This nonnegative scalar combines the two normal components and counts shear
once. Squaring removes their signs, so deformation direction is lost. The
original flow components accompany it as separate channels.

For comparison, the chapter's prose notes the extra shear contribution in the
Frobenius norm, in which both symmetric off-diagonal entries enter. Chapter 3
is where the thesis actually prints that four-term convention, so under this
guide's no-repetition rule it is given once, at
Equation~\eqref{fg:eq:frobenius}, together with the argument for departing
from it. The difference changes shear's relative weight and is not merely a
single global rescaling.

\textbf{Source.} The tensor components come from Shreve
et al.~\cite{shreve2011}; the exact three-term scalar is verified in
\path{Stage1_DataPipeline/step2_extraction/flow_strain_extractor.py}.
It is not claimed to be Shreve's identical aggregation rule. The Frobenius
expression follows from the standard matrix-norm definition.

\subsubsection{Tensor shape and channel normalisation}
\label{fg:tensor}
\textbf{Thesis location:} 2.4.4.

The channel order $[u,v,\mathrm{strain}]$ produces shape $(3,32,224,224)$:
three channels, 32 adjacent-frame fields and a $224\times224$ spatial grid.
The 32 fields come from 33 selected frames. These are dimensions, not a
formula from a paper, and the spatial grid covers the resized full frame.

\label{fg:input-normalisation}
The chapter describes two transformations in prose. Their code-equivalent
formulas are added here for clarity:
\begin{equation}
    z_i=\frac{x_i-x_{\min}}{x_{\max}-x_{\min}+10^{-8}},
    \qquad
    \widehat a_i=\frac{a_i-\mu_a}{s_a+10^{-6}}.
    \label{fg:eq:input-normalisation}
\end{equation}
Extrema, mean and standard deviation are computed separately for each channel
over its $T\times H\times W$ values. Extraction maps into $[0,1]$.
In the second expression, $a$ is the stored channel after any training
augmentation; evaluation applies no augmentation. Loading standardises $a$,
using its own mean $\mu_a$ and standard deviation $s_a$. The loader's standard
deviation uses the default sample correction. The small constants prevent
division by zero, and a constant channel need not have unit variance. The
standardisation stabiliser is added to the standard deviation, unlike the
inside-square-root stabiliser in batch normalisation below.

\textbf{Source.} Project implementation in
\texttt{flow\_strain\_extractor.py} and \texttt{Ablation\_Study/dataset.py}.
These are standard normalisation operations with code-specific scope and
constants; no originating research paper is asserted for their exact combination.

\subsection{Neural-network fundamentals: thesis section 2.5}

\subsubsection{A computational unit}
\label{fg:unit}
\textbf{Thesis location:} 2.5.1.
\begin{equation}
    a=\phi(w^{\mathsf T}x+b).
    \label{fg:eq:unit}
\end{equation}
The input vector $x$ is combined with learned weights $w$ by a dot product;
the bias $b$ shifts that result; the activation $\phi$ produces output $a$.
The activation's shape is fixed while the surrounding weights are learned.

\textbf{Source.} This is the standard affine-unit definition used in the
background chapter, not an invention attributed to a particular supplied paper.
OFF-ApexNet~\cite{liong2019a}, PDF page 6, equation (12), provides a related
convolutional weighted-sum, bias and activation formulation; it is contextual
support rather than the literal origin of this compact vector notation.

\subsubsection{Rectified linear activation}
\label{fg:relu}
\textbf{Thesis location:} 2.5.2.
\begin{equation}
    \operatorname{ReLU}(x)=\max(0,x).
    \label{fg:eq:relu}
\end{equation}
Negative inputs become zero and positive inputs pass through unchanged. This
nonlinearity prevents a stack of affine layers from reducing to one affine
mapping. The source chapter names GELU as well, but gives no GELU equation.

\textbf{Source.} OFF-ApexNet~\cite{liong2019a}, PDF page 6, equation (13),
explicitly presents this formula. It is a supplied source for the definition,
not the original ReLU paper.

\subsubsection{Convolutional parameter count}
\label{fg:convparams}
\textbf{Thesis location:} 2.5.4.
% SOURCE-DISPLAY-08
\begin{equation}
    N_{\mathrm{params}}=C_{\mathrm{out}}
       (C_{\mathrm{in}}k_Hk_W+1).
    \label{fg:eq:convparams}
\end{equation}
For an ungrouped convolution, each of the $C_{\mathrm{out}}$ output channels has $C_{\mathrm{in}}k_Hk_W$
weights and one bias. Here $C_{\mathrm{in}}$ counts input channels and
$k_H,k_W$ are kernel height and width. Weight sharing makes this count
independent of image size, although computation and activation memory still
grow with spatial resolution. If bias is disabled, the $+1$ is omitted; the
actual stem's convolution layers disable it.

\textbf{Source.} Direct parameter counting, not a borrowed empirical result
or a formula needing an invented source citation. The bias setting is verified
in \texttt{Ablation\_Study/models.py}.

\subsubsection{Pooling notation}
\label{fg:pool}
\textbf{Thesis location:} 2.5.5.

The inline $2\times2$ denotes the spatial pooling window. With stride two and
compatible dimensions, it halves height and width. It is a shape setting;
the shortened chapter does not display a pooling equation. Neither maximum
nor average pooling learns weights. A later projection is a separate learned
operation. This interpretation follows the implemented layer definitions,
rather than attributing the exact settings to a paper.

\subsection{Network blocks: thesis section 2.6}

\subsubsection{Three-dimensional tensor and kernel notation}
\label{fg:conv3d}
\textbf{Thesis location:} 2.6.1.

The tensor axes $(B,C,T,H,W)$ mean batch, channels, time, height and width.
The kernel $(k_T,k_H,k_W)$ directly mixes neighbouring times only when
$k_T>1$. The stem uses $(1,3,3)$ kernels, $(0,1,1)$ padding and $(1,2,2)$
pooling. Thus these kernels preserve time length and perform spatial
convolution. Whole-clip statistics in normalisation or SimAM can still create
dependence across frames. All these tuples are implementation settings,
verified in \texttt{Ablation\_Study/models.py}.

\subsubsection{Batch normalisation}
\label{fg:bn}
\textbf{Thesis location:} 2.6.2.
% SOURCE-DISPLAY-09
\begin{equation}
    \widehat x=\frac{x-\mu_c}{\sqrt{\sigma_c^2+\epsilon}},
    \qquad y=\gamma_c\widehat x+\beta_c.
    \label{fg:eq:bn}
\end{equation}
For channel $c$, subtract mean $\mu_c$, divide by the stabilised standard
deviation and apply learned scale $\gamma_c$ and shift $\beta_c$. Training
statistics for BatchNorm3d span $(B,T,H,W)$; evaluation uses running estimates.
These activation statistics differ from per-sample input normalisation.

\textbf{Source.} A standard layer definition, consistent with the
\texttt{BatchNorm3d} layers in \texttt{Ablation\_Study/models.py}.
The source chapter gives no originating batch-normalisation paper. No such
paper is claimed here merely because another reviewed model uses the layer.

\subsubsection{SimAM statistics}
\label{fg:simamstats}
\textbf{Thesis location:} 2.6.3.
% SOURCE-DISPLAY-10
\begin{equation}
    \mu=\frac1M\sum_i x_i,\qquad
    v=\frac1{M-1}\sum_i(x_i-\mu)^2.
    \label{fg:eq:simamstats}
\end{equation}
$x_i$ is an activation at one position in a channel, $M$ counts positions,
$\mu$ is their mean and $v$ their variance estimate. Here $M=THW$, so each
sample's statistics span the full clip; the original image setting uses $HW$.
Both sums run over all positions $i=1,\ldots,M$, and the sample variance
requires $M>1$, as satisfied by this feature volume.

\textbf{Source.} Yang et al.~\cite{yang2021}, PDF pages 4--5,
equation (5) and Figure 3. The printed variance definition divides by $M$,
but Figure 3's reference code divides by $M-1$. The thesis follows the code
and extends its context from images to clips. It must not be described as
an unchanged copy of the printed variance equation.

\subsubsection{SimAM energy and multiplicative gate}
\label{fg:simamgate}
% SOURCE-DISPLAY-11
\begin{equation}
\begin{aligned}
    E_i^{-1}&=\frac{(x_i-\mu)^2}{4(v+\lambda)}+\frac12,\\
    \widetilde x_i&=x_i\operatorname{sigmoid}(E_i^{-1}).
\end{aligned}
    \label{fg:eq:simamgate}
\end{equation}
The squared deviation measures distinctiveness from the channel mean.
Dividing by $v+\lambda$ puts it in the context of channel variation;
$\lambda=10^{-4}$ prevents an unstable division for low variance.
$E_i^{-1}$ is reciprocal energy. The sigmoid turns it into a bounded gate
that scales $x_i$, preserving score order but compressing its values.
No learned gate parameters are added, although calculating the gate costs
time and memory. Distinctive activation is not automatically useful motion.

\textbf{Source.} Yang et al.~\cite{yang2021}, equations (5)--(6) and
Figure 3. Taking the reciprocal of the paper's minimum-energy expression
produces the displayed score; the variance convention and clip context are
the implementation adaptations described above.

\subsubsection{Scaled dot-product attention}
\label{fg:attention}
\textbf{Thesis location:} 2.6.4.
% SOURCE-DISPLAY-12
\begin{equation}
    \operatorname{Attention}(Q,K,V)
      =\operatorname{softmax}\left(\frac{QK^{\mathsf T}}{\sqrt{d_k}}\right)V.
    \label{fg:eq:attention}
\end{equation}
$Q,K,V$ are query, key and value matrices derived from input tokens.
$QK^{\mathsf T}$ scores each query against the keys, and $d_k$ is the key
dimension. Dividing by $\sqrt{d_k}$ controls the growth of dot products;
row-wise softmax produces nonnegative weights summing to one. Multiplication
by $V$ makes each output a weighted combination of values. This is the
single-head operation, not the complete transformer block.

\textbf{Source.} Dosovitskiy et al.~\cite{dosovitskiy2021}, supplied PDF
page 13, Appendix A, equations (5)--(8), explicitly define this attention
operation with per-head dimension $D_h$, equivalent to $d_k$ here. This is a
verified supplied presentation of the standard transformer operation, not
its original invention. The thesis uses fixed positions; the equation alone
contains no position information and does not establish temporal-order sensitivity.

\subsubsection{The learned no-CNN projection}
\label{fg:projection}
\textbf{Thesis location:} 2.6.5.
\begin{equation}
    N_{\mathrm{projection}}=48\times96+96=4{,}704.
    \label{fg:eq:projection}
\end{equation}
Pooling three channels to $4\times4$ gives 48 features per frame. The
$48\rightarrow96$ linear layer then learns 96 sets of 48 weights and 96
biases. Removing the CNN therefore does not remove all learned input
processing. This is direct arithmetic from \texttt{RawPatchEmbedding} in
\texttt{Ablation\_Study/models.py}, not a count imported from a paper.

\subsection{Training under scarcity and skew: thesis section 2.7}

\subsubsection{Cross-entropy for a hard target}
\label{fg:ce}
\textbf{Thesis location:} 2.7.1.
\begin{equation}
    L_{\mathrm{CE}}=-\log p_t.
    \label{fg:eq:ce}
\end{equation}
Here $t$ indexes the true class, not time, and $p_t$ is its predicted
probability. A confident correct prediction has small loss; a small true-class
probability has large loss. The logarithm is natural. Under ordinary sampling,
frequent classes contribute more examples, but class frequency alone does not
determine each example's gradient.

\textbf{Source.} Standard negative log-likelihood. Zhang
et al.~\cite{zhang2022}, supplied PDF page 6, section 3.4, equation (15),
present multiclass cross-entropy; reducing its class sum for a one-hot target
gives this per-example expression. Zhao et al.~\cite{zhaoS2021} present the
focal extension below, whose unweighted zero-exponent case also reduces to it.

\subsubsection{Focal loss}
\label{fg:focal}
\textbf{Thesis location:} 2.7.2.
% SOURCE-DISPLAY-13
\begin{equation}
    \operatorname{FL}(p_t)=-\alpha_t(1-p_t)^\gamma\log p_t.
    \label{fg:eq:focal}
\end{equation}
For $\gamma>0$, the factor $(1-p_t)^\gamma$ suppresses easy, confident examples
relative to hard ones. $\alpha_t$ is an optional class weight, serving a
different purpose from difficulty weighting. The study uses $\gamma=2$.
For example, the modulating factor is 0.01 at $p_t=0.9$ and 0.81 at $p_t=0.1$.
These are modulation factors, not complete loss values.

\textbf{Source.} Zhao et al.~\cite{zhaoS2021}, PDF page 7, section 3.2.2.2,
equation (10), apply focal loss to micro-expression recognition and report
$\gamma=2$. This is an application source, not a claim that Zhao et al.
originated focal loss. The thesis additionally applies label smoothing,
which must be distinguished from this unsmoothed expression.

\subsubsection{Inverse-frequency sampling}
\label{fg:sampling}
\textbf{Thesis locations:} 2.7.3--2.7.4.
\begin{equation}
    w_i=\frac1{n_{c_i}}.
    \label{fg:eq:sampling}
\end{equation}
$n_{c_i}$ is the number of training examples in example $i$'s class, counted
within the current fold. Each represented class has total weight
$n_c(1/n_c)=1$, so normalised sampling probabilities give those classes
equal mass. This is an expectation over replacement draws; individual
batches need not be balanced and an epoch can repeat or omit clips.
The held-out loader is not rebalanced. When this sampler is active, the
implementation disables loss-side class weights to avoid applying two
frequency corrections at once.

\textbf{Source.} The weights and replacement rule are project choices in
\texttt{tools/run\_ablation\_experiments.py}. The equal-mass result is their
algebraic consequence, not an independently sourced research finding.

\subsubsection{Smoothed targets and smoothed cross-entropy}
\label{fg:smoothing}
\textbf{Thesis location:} 2.7.5.
\begin{equation}
    q_c=(1-\epsilon)\mathbf{1}[c=t]+\epsilon/C.
    \label{fg:eq:smoothing-target}
\end{equation}
There are $C$ classes. The indicator is one for the true class and zero
otherwise; $\epsilon$ redistributes some target mass uniformly. With
$C=3$ and $\epsilon=0.05$, the true-class target is approximately 0.9667 and
each other target approximately 0.0167. The targets sum to one.
% SOURCE-DISPLAY-14
\begin{equation}
    L_{\mathrm{smooth}}
      =-(1-\epsilon)\log p_t
        -\frac{\epsilon}{C}\sum_{c=1}^{C}\log p_c.
    \label{fg:eq:smoothing-loss}
\end{equation}
This is cross-entropy against the smoothed targets, obtained by substituting
them into the sum over class log-probabilities. It discourages extreme
confidence rather than preferentially weighting a minority class.

The chapter also describes the implemented combination in prose. Written
out as an explanatory expansion, it is
\begin{equation}
    L_{\mathrm{implemented}}=\alpha_t(1-p_t)^\gamma L_{\mathrm{smooth}},
    \label{fg:eq:combined-loss}
\end{equation}
with $\alpha_t$ effectively one when class weighting is disabled. The focal
factor still uses the unsmoothed hard-target probability $p_t$. This is the
per-example loss; the implemented default reduction averages it over the batch.

\textbf{Source.} Exact combination and smoothing strength verified in
\texttt{Ablation\_Study/losses.py}. The target distribution is a standard
definition and the loss is its algebraic expansion. No supplied paper is
claimed to originate this exact focal-plus-smoothing implementation; Zhao
et al.~\cite{zhaoS2021} support the focal component only.

\textbf{Remaining training prose.} Section 2.7.6 describes horizontal flipping
and negating the horizontal-flow channel $u$; it gives no further displayed
formula. Section 2.7.7 names AdamW, warmup and cosine annealing but writes no
update or schedule equation. Their full equations have not been invented as
extractions from this chapter.

\subsection{Evaluation: thesis section 2.8}

\subsubsection{Precision, recall and class F1}
\label{fg:metrics}
\textbf{Thesis location:} 2.8.1.
% SOURCE-DISPLAY-15
\begin{equation}
\begin{aligned}
    P_c&=\frac{TP_c}{TP_c+FP_c},\qquad
    R_c=\frac{TP_c}{TP_c+FN_c},\\
    F1_c&=\frac{2TP_c}{2TP_c+FP_c+FN_c}.
\end{aligned}
    \label{fg:eq:metrics}
\end{equation}
For class $c$, $TP_c$ counts correct predictions of that class, $FP_c$ counts
other examples incorrectly predicted as that class, and $FN_c$ counts its
examples assigned elsewhere. Precision asks how many positive predictions
are correct; recall asks how many actual positives were found. F1 balances
both through their harmonic mean. In the confusion matrix, entry $(i,j)$
counts true class $i$ predicted as class $j$.

The project sets undefined class scores to zero. Changing a single prediction
does not have a fixed effect on F1 because both its numerator and denominator
depend on the existing confusion counts.

\textbf{Source.} OFF-ApexNet~\cite{liong2019a}, supplied PDF page 9,
equations (20)--(22), explicitly gives recall, precision and their harmonic-mean
F1; substituting the first two into the third gives the count-based F1 above.
See et al.~\cite{see2019}, PDF page 3, equation (3), give the direct count-based
F1 expression. These support the standard definitions, not a claim that those
papers originated precision or recall.

\subsubsection{Macro F1 and micro F1}
\label{fg:macro}
\textbf{Thesis location:} 2.8.2.
% SOURCE-DISPLAY-16
\begin{equation}
    F1_{\mathrm{macro}}=\frac1C\sum_{c=1}^{C}F1_c.
    \label{fg:eq:macro}
\end{equation}
Every class contributes equally, regardless of support. Thus Negative does
not receive greater explicit weight merely because it supplies more clips.
The chapter describes micro aggregation in prose; its explanatory expansion is
\begin{equation}
\begin{aligned}
    F1_{\mathrm{micro}}&=
    \frac{2\sum_c TP_c}{2\sum_c TP_c+\sum_c FP_c+\sum_c FN_c}\\
    &=\frac{N_{\mathrm{correct}}}{N}
    \quad\text{for single-label multiclass prediction}.
\end{aligned}
    \label{fg:eq:micro}
\end{equation}
Here $N$ is the number of evaluated clips and $N_{\mathrm{correct}}$ is the
number classified correctly. Each incorrect classification contributes one false positive and one false
negative across the classes. Micro F1 therefore equals accuracy in this task;
the identity does not hold for every possible classification setting.

\textbf{Source.} Macro F1 and the pooled class-average construction are
supported by See et al.~\cite{see2019}, PDF page 3. The micro-F1 equality is
an algebraic consequence of single-label prediction, not a new experimental
result from that paper.

\subsubsection{Pooling folds before calculating F1}
\label{fg:fold-pooling}
\textbf{Thesis locations:} 2.8.3--2.8.5.

The chapter states the following distinction in prose. Let $A^{(k)}$ be the
confusion matrix of held-out subject $k$ and let $K$ be the number of folds.
An explicit explanatory form is
\begin{equation}
\begin{aligned}
    A_{\mathrm{pooled}}&=\sum_{k=1}^{K}A^{(k)},\\
    F1_{\mathrm{macro}}(A_{\mathrm{pooled}})
       &\ne \frac1K\sum_{k=1}^{K}F1_{\mathrm{macro}}(A^{(k)})
       \quad\text{in general}.
\end{aligned}
    \label{fg:eq:pooling}
\end{equation}
Compute class F1 after summing counts for the primary pooled score; do not
substitute the average of fold scores. F1 is nonlinear and folds can lack
classes. Equality is possible in special cases, so the inequality is not
asserted for every dataset. Subject-disjoint folds also do not undo selection
bias when the same held-out labels select a checkpoint and score it.

\textbf{Source.} See et al.~\cite{see2019}, PDF page 3, explicitly accumulate
confusion counts over LOSO folds before computing Unweighted F1. Its printed
equation (4) appears to omit the summation; the guide follows its explanatory
prose and the mathematically complete class average. The matrix notation above
is a formalisation added by this guide. The thesis uses $K=25$; its stored
\texttt{macro\_f1} field is the mean of fold scores, whereas its reported
primary score is reconstructed from pooled counts.

\subsection{Chapter 2 coverage boundary}
All sixteen displayed equation blocks in the shortened Chapter 2 are included
above, together with its inline formulas, tuple notation and numerical ranges,
with the single exception noted at \autoref{fg:magnitude}: the Frobenius
comparison is printed once, in Chapter 3, where the thesis states it as a
published convention. The concluding section 2.9 has no further formulas.
Topics mentioned without an equation have not been expanded into a separate
textbook treatment.

\section{Chapter 3: Literature Review}
\label{fg:chapter3}

\textbf{Source file:} \texttt{contents/chapter3\_shortened.tex}. The chapter
contains \textbf{seven displayed equation blocks}. Four of them restate
relations already given in Chapter~2 --- the strain magnitude, the SimAM
statistics, the SimAM gate and focal loss --- because Chapter~2 develops the
mechanism and Chapter~3 reports the literature's version of it. Following the
no-repetition rule of this guide, those are not printed a second time. Each is
identified below by its Chapter~2 equation number, with only the material
Chapter~3 genuinely adds: the published convention it departs from, the
divergence, and the argument for it.

\subsection{Evaluation protocol: thesis section 3.1}

\subsubsection{Unweighted F1 and the MEGC metric convention}
\label{fg:uf1}
\textbf{Thesis location:} 3.1.6, \texttt{sec:megc-metrics}.

\textbf{No new formula.} Chapter~3 prints
$F1_c = 2TP_c/(2TP_c + FP_c + FN_c)$ and
$\text{UF1} = (1/C)\sum_c F1_c$. These are Equations
\eqref{fg:eq:metrics} and \eqref{fg:eq:macro} exactly, under the name the
Second Facial Micro-Expression Grand Challenge gives them. Unweighted F1
\emph{is} macro-averaged F1; the guide does not restate the algebra.

What Chapter~3 adds is the \textbf{construction}, which the algebra does not
show. UF1 is a \textbf{pooled} quantity: $TP_c$, $FP_c$ and $FN_c$ accumulate
across all $k$ folds \emph{before} any $F1_c$ is computed. That is the
distinction formalised at Equation~\eqref{fg:eq:pooling}, and it separates UF1
from the average of per-fold macro F1 scores, which is a different estimator
with a different expectation. The chapter also fixes the protocol's scope: the
Composite Database Evaluation merges SMIC, CASME~II and SAMM into 442
samples --- Negative 250, Positive 109, Surprise 83 --- of which CASME~II
contributes 145, with leave-one-subject-out run across the composite. Plain
accuracy is rejected explicitly on grounds of class imbalance.

This thesis adopts the metric definitions while declining the composite
training regime. Its primary score, pooled macro F1, is identical in
construction to UF1.

\textbf{Source.} See et al.~\cite{see2019}. The per-class F1 and macro-average
attributions are those already given for Equations~\eqref{fg:eq:metrics}
and~\eqref{fg:eq:macro}.

\subsection{Optical strain: thesis section 3.4}

\subsubsection{The four-term published magnitude}
\label{fg:frobenius}
\textbf{Thesis location:} 3.4.2, \texttt{sec:strain-tensor}.

This is the convention the thesis departs from, and Chapter~3 is where the
thesis prints it. Chapter~2 alludes to it in prose only
(\autoref{fg:magnitude}).
% SOURCE-DISPLAY-CH3-01
\begin{equation}
    \varepsilon_{\mathrm{F}}
    =\sqrt{\varepsilon_{xx}^{2}+\varepsilon_{yy}^{2}
        +\varepsilon_{xy}^{2}+\varepsilon_{yx}^{2}}
    =\sqrt{\varepsilon_{xx}^{2}+\varepsilon_{yy}^{2}
        +2\varepsilon_{xy}^{2}}.
    \label{fg:eq:frobenius}
\end{equation}
The components are those of Equations~\eqref{fg:eq:strain-normal}
and~\eqref{fg:eq:strain-shear}. This is the Frobenius norm of the symmetric
tensor: because the tensor is symmetric, both off-diagonal entries enter the
sum, and the second equality follows. The scalar describes deformation
intensity and discards its sign and direction. Supplying it together with $u$
and $v$ retains the displacement fields from which those spatial relations
were calculated.

Discretisation is a modelling choice rather than a neutral detail, because it
sets the scale at which local variation is measured: Shreve et
al.~\cite{shreve2011} use central differences with offsets of approximately
two to three pixels, whereas Liong et al.~\cite{liong2014a} use one-pixel
offsets.

\textbf{Source.} The Liong strain series~\cite{liong2014a,liong2014b,liong2016}
uses this four-component magnitude as its descriptor. Shreve et
al.~\cite{shreve2011} establish the tensor itself, cited at
\autoref{fg:strain}; the four-term aggregation is the Liong convention.

\subsubsection{The implemented magnitude and the shear-weighting divergence}
\label{fg:strain-divergence}
\textbf{Thesis location:} 3.4.3, \texttt{sec:strain-divergence}.

\textbf{The implemented scalar is not repeated here.} Chapter~3 writes it as
$\varepsilon_{\mathrm{os}}$ and Chapter~2 as $\varepsilon_{\mathrm{mag}}$;
they are the same three-term expression, Equation~\eqref{fg:eq:strain-magnitude}.
Chapter~3 restates the shear component alongside it, which is
Equation~\eqref{fg:eq:strain-shear}.

What Chapter~3 adds is the divergence and its justification. The implemented
form \textbf{counts shear once}; the published convention of
Equation~\eqref{fg:eq:frobenius} counts it \textbf{twice}. The four-term form
therefore weights pure shear by $\sqrt{2}$ relative to the implemented one.
Independent min--max normalisation removes a uniform scale factor but does not
generally remove this difference, because normal and shear contributions vary
across pixels and across time. The two conventions were never compared
experimentally, so their practical equivalence is not assumed.

STSTNet~\cite{liong2019b} cannot establish a precedent for the three-term
form, because its printed magnitude equation is defective in two independent
ways: it repeats $\partial u$ in the mixed derivative where the symmetric
tensor calls for $\partial v/\partial x$ alongside $\partial u/\partial y$,
and it places the factor $\tfrac{1}{2}$ \emph{outside} the squared derivative
sum. Correcting the repeated derivative while retaining that placement gives
% SOURCE-DISPLAY-CH3-02
\begin{equation}
    \tfrac{1}{2}
       \left(\frac{\partial u}{\partial y}
            +\frac{\partial v}{\partial x}\right)^2
       =2\varepsilon_{xy}^{2},
    \label{fg:eq:shear-identity}
\end{equation}
which agrees with the \emph{four}-term convention, not the three-term one.
Moving the factor inside the square would instead yield the implemented form,
but that is a further change to the printed expression. The intended
implementation cannot be settled from the published equation alone.

The thesis therefore declares a departure from the unambiguous four-term
convention rather than claiming to reproduce STSTNet's magnitude. This is a
divergence from the whole corpus, not a following of STSTNet, and it is
recorded as row~4 of the divergence ledger. \textbf{Neither expression is a
typographical error to be corrected: the difference between them is the
finding.}

\textbf{Source.} No source. This is a declared divergence. Spatial
derivatives come from NumPy's \texttt{gradient} at one-pixel spacing, using
central differences in the array interior --- spacing that follows Liong et
al.~\cite{liong2014a} even though the scalar reduction does not.

\subsubsection{Flow magnitude}
\label{fg:flow-magnitude}
\textbf{Thesis location:} 3.3.2, \texttt{sec:brightness-constancy}.

The chapter's inline expression, promoted here, is
\begin{equation}
    \|(u,v)\|=\sqrt{u^2+v^2}.
    \label{fg:eq:flow-magnitude}
\end{equation}
This is the scalar magnitude of the flow vector at a pixel, discarding its
direction. It is included because the thesis \emph{declines} it: the two
components are retained as separate input channels, so direction of motion
survives into the network. The strain channel of
Equation~\eqref{fg:eq:strain-magnitude} is the only scalar reduction the
pipeline performs.

\textbf{Source.} Zhao et al.~\cite{zhaoS2021} describe this magnitude
alongside the component fields $u(x,y)$ and $v(x,y)$. The decision to keep the
components instead is a project choice.

\subsection{Attention and temporal normalisation: thesis sections 3.6 and 3.5}

\subsubsection{SimAM statistics and gate as restated in Chapter 3}
\label{fg:simam-ch3}
\textbf{Thesis location:} 3.6.3, \texttt{sec:simam-energy}.

\textbf{No new formula.} Chapter~3 writes the module with an intermediate
squared deviation $d_i=(x_i-\mu)^2$ and shared variance $s^2$, so that the
score reads $E_i^{-1}=d_i/(4(s^2+\lambda))+\tfrac12$. Substituting $d_i$ and
renaming $s^2$ to $v$ gives Equations~\eqref{fg:eq:simamstats}
and~\eqref{fg:eq:simamgate} unchanged. The two chapters differ in notation
only.

What Chapter~3 adds is the provenance of the $M-1$ denominator and the status
of $\lambda$. The shared variance printed with Yang et al.'s equation~(5)
divides by $M$, whereas the reference code in their Figure~3 divides by
$M-1$; the thesis follows the \emph{code}, applying Bessel's correction over
$n = D \cdot H \cdot W - 1$. The $+\tfrac12$ follows from the
reciprocal-energy expression rather than being an arbitrary offset, and the
sigmoid preserves score ordering while bounding the weights --- it does not
produce a hard mask.

The value $\lambda=10^{-4}$ is inherited, not derived. Yang et al. sweep
$\lambda$ from $10^{-1}$ to $10^{-6}$ and select $10^{-4}$ for CIFAR,
repeating the selection separately for ImageNet and choosing $0.1$. It is a
dataset-selected operating point, not a setting the formula justifies.

Chapter~3 also fixes what the module's ``three-dimensional weights'' mean: the
original produces a separate weight at each channel--height--width position,
so they describe $C\times H\times W$ and not a temporal volume. Extending the
context to $M=THW$ is an additional design decision, recorded as row~10 of the
divergence ledger. Whole-clip statistics do not guarantee that apex frames
receive the highest weights.

\textbf{Source.} Yang et al.~\cite{yang2021}, as for
Equations~\eqref{fg:eq:simamstats} and~\eqref{fg:eq:simamgate}.

\subsubsection{Effective frame rate after subsampling}
\label{fg:effective-rate}
\textbf{Thesis location:} 3.5.6, \texttt{sec:tim-implications}.

The chapter's inline expression, promoted here, is
\begin{equation}
    f_{\mathrm{eff}}=\frac{33}{N}\,f_s=\frac{6600}{N}\ \text{fps}
    \qquad (f_s=200).
    \label{fg:eq:effective-rate}
\end{equation}
Sampling 33 frames from a clip of $N$ recorded frames spreads them across the
clip's full duration, so the images reaching the model are separated by more
than the recording interval. At the median clip length $N=66$ this gives
\textbf{100 fps} --- SMIC's recording rate, and half what CASME~II was built to
provide. At the longest clip, $N=126$, it gives roughly \textbf{52 fps}, below
the original CASME database's 60 fps~\cite{yan2014}. Eighty-one of the 156
clips are downsampled by a factor of two or more.

This is the cost of fixed-length temporal normalisation, expressed in the unit
the capture rate was chosen in (Equation~\eqref{fg:eq:capture}). No reviewed
paper reports it.

\textbf{Source.} This thesis's own arithmetic over the working manifest. The
200 fps capture rate is Yan et al.~\cite{yan2014}.

\subsubsection{Realised magnification pass-band}
\label{fg:evm-interval}
\textbf{Thesis location:} 3.5.6, \texttt{sec:tim-implications}.

The chapter's inline expression, promoted here, is
\begin{equation}
    \Delta t_{\mathrm{true}}=\frac{N}{200\times 32}\ \text{seconds}.
    \label{fg:eq:evm-interval}
\end{equation}
Magnification runs on the 33 sampled frames at the \emph{nominal} 200 fps, but
those frames span the clip's full duration, so the true interval between
consecutive sampled frames is the above rather than $1/200$. The filter's
frequency axis is therefore scaled by $N/32$ relative to what the
implementation assumes.

The consequence is that a band specified as 5--25 Hz corresponds to roughly
2.5--12.5 Hz physically at the median, with the discrepancy growing with clip
length. The realised pass-band is thus \textbf{clip-dependent} --- neither the
5--25 Hz derived from the duration rule of Bai et al.~\cite{bai2021} nor
constant across the corpus. This is a declared divergence; magnifying before
subsampling is declared as further work.

\textbf{Source.} This thesis's own arithmetic. The band specification derives
from Bai et al.~\cite{bai2021}; the discrepancy is a project finding.

\subsection{Class imbalance: thesis section 3.9}

\subsubsection{Focal loss as restated in Chapter 3}
\label{fg:focal-ch3}
\textbf{Thesis location:} 3.9.2, \texttt{sec:imbalance-focal}.

\textbf{No new formula.} Chapter~3 prints
$\text{FL}(p_t) = -\alpha_t(1-p_t)^{\gamma}\log p_t$, which is
Equation~\eqref{fg:eq:focal} exactly.

What Chapter~3 adds is the division of labour between the two factors, and it
matters because conflating them is where this project went wrong. Zhao et
al.~\cite{zhaoS2021} describe $\gamma$ as ``the balance factor for loss'' and
$\alpha_t$ as ``the weight balance factor for samples''. The modulating factor
$(1-p_t)^\gamma$ re-weights by \textbf{difficulty} and does not inspect class
frequency at all, so a difficult example from any class retains a large
contribution --- though under skew the examples it up-weights are
disproportionately from the minority classes. The separate $\alpha_t$
re-weights by \textbf{class frequency}. Zhao et al. treat $\gamma$ as a
constant, reporting ``$\gamma$ is set as 2 in practice'', and $\alpha$ as
dataset-specific, ``to set by cross validation''.

This thesis assigns frequency correction to the sampler of
Equation~\eqref{fg:eq:sampling} instead, and a runtime guard disables the
loss-side class weights whenever balanced sampling is active, leaving focal
difficulty weighting in place. The implemented composition is
Equation~\eqref{fg:eq:combined-loss}.

\textbf{Source.} Zhao et al.~\cite{zhaoS2021} bring focal loss into
micro-expression recognition, adopting the formulation developed for one-stage
object detection where the imbalance is between foreground and background
regions. This is an application source, not a claim that they originated focal
loss; the originating detection paper is not in the thesis's corpus and is not
cited here.


\section{Chapter 4: Methodology}
\label{fg:chapter4}

\textbf{Source file:} \texttt{contents/chapter4\_shortened.tex}. The chapter
contains \textbf{two displayed equation blocks}, both unique to it, together
with several inline counts. It is the reproduction chapter, so parameter
values, configuration names, tensor shapes and seeds are stated exactly.

\subsection{The configuration space: thesis section 4.2}

\subsubsection{The factorial design space}
\label{fg:factorial}
\textbf{Thesis location:} 4.2.2, \texttt{sec:twelve-configs}.

The chapter's inline expression, promoted here, is
\begin{equation}
    N_{\mathrm{cells}}=2^4=16,
    \qquad N_{\mathrm{valid}}=16-4=12.
    \label{fg:eq:factorial}
\end{equation}
Four binary flags give sixteen cells, constructed directly as the Cartesian
product of the flag states
(\verb|itertools.product([False, True], repeat=4)|). Each candidate then
passes \verb|AblationConfig.is_valid()|, which rejects any configuration with
\texttt{use\_simam=True} while \texttt{use\_cnn=False}. Exactly four cells
carry that combination and all four are rejected, leaving \textbf{twelve}
trained and evaluated.

The rejection is architectural, not empirical: SimAM re-weights activations in
a convolutional feature map, so without a convolutional stem there is no map
for it to act on. Such a cell is not a weaker pipeline but a specification
that does not correspond to a well-formed model.

\textbf{Source.} Direct arithmetic over the implementation's flag set. No
paper is claimed for the count.

\subsection{The transformer: thesis section 4.5}

\subsubsection{Sinusoidal positional encoding}
\label{fg:posenc}
\textbf{Thesis location:} 4.5.3, \texttt{sec:transformer-config}.
% SOURCE-DISPLAY-CH4-01
\begin{equation}
	\mathrm{PE}(p, 2i) = \sin\!\left(\frac{p}{10000^{2i/96}}\right), \qquad \mathrm{PE}(p, 2i+1) = \cos\!\left(\frac{p}{10000^{2i/96}}\right).
    \label{fg:eq:posenc}
\end{equation}
$p$ is the position in the sequence, one of the 32 time steps; $i$ indexes
feature pairs, so $2i$ and $2i+1$ are the even and odd feature positions; 96
is $d_{\text{model}}$, appearing as a literal because the configuration fixes
it. Each feature pair oscillates at a frequency set by the $10000^{2i/96}$
denominator, geometrically spaced across the embedding, so every position
receives a distinct vector and positions at a fixed offset stand in a
consistent linear relation.

The encoding exists because attention needs it. Equation~\eqref{fg:eq:attention}
is permutation-equivariant on its own: permuting the input tokens permutes the
outputs correspondingly, so a shuffled clip would be indistinguishable from an
ordered one. The positional encoding is what makes the temporal axis
meaningful.

It is a fixed buffer, not a learned parameter, and contributes nothing to any
parameter count. Two details are easily missed: the positional-encoding stage
applies its own \texttt{Dropout(p=0.1)} \emph{after} the encoding is added,
which is a separate application from the encoder layers' dropout at the same
rate; and although the sinusoids could be evaluated at arbitrary positions,
the implementation uses a finite precomputed buffer, so generalisation to
unseen sequence lengths is neither exercised nor guaranteed.

\textbf{Source.} \textbf{No corpus source.} The encoding is the standard
transformer form, which Chapter~4 states without attribution. The thesis
reaches the architecture through Zhang et al.'s SLSTT~\cite{zhang2022} and its
Vision Transformer foundation~\cite{dosovitskiy2021}; note that ViT itself
uses \emph{learned} positional embeddings, not these fixed sinusoids, so
neither supplied paper is claimed to present this equation.

\subsection{Parameter budget: thesis section 4.6}

\subsubsection{Component counts and configuration totals}
\label{fg:param-totals}
\textbf{Thesis location:} 4.6.1, \texttt{sec:parameter-counts}.

Four components combine to give any configuration's total: the convolutional
backbone (\texttt{ThreeStreamCNNBackbone}, all three streams)
\textbf{14,544}; the transformer encoder (\texttt{SLSTTTransformer})
\textbf{348,736}; \texttt{RawPatchEmbedding}, the no-CNN fallback,
\textbf{4,704}, derived at Equation~\eqref{fg:eq:projection}; and the
classifier head \textbf{483}. Their sums, written out here as an explanatory
expansion of the chapter's table, are
\begin{equation}
\begin{aligned}
    4{,}704+483&=5{,}187, &\qquad 14{,}544+483&=15{,}027,\\
    4{,}704+348{,}736+483&=353{,}923, &\qquad
    14{,}544+348{,}736+483&=363{,}763.
\end{aligned}
    \label{fg:eq:param-totals}
\end{equation}
The four rows are, in order: no CNN and no transformer; CNN without
transformer; transformer without CNN; and both. Every configuration ends in
the same classifier head, so 483 appears in each.

Neither SimAM nor EVM appears as a term. SimAM contributes zero parameters by
construction --- no weight or bias is learned anywhere in
Equation~\eqref{fg:eq:simamgate} --- and magnification is a preprocessing step
rather than a network component. Of the transformer's 348,736, exactly 192
belong to the final \texttt{LayerNorm(96)} applied after its four encoder
layers.

In a full configuration the transformer accounts for roughly \textbf{96\,\%}
of all parameters (348,736 of 363,763). That is a fact about the parameter
budget, not a claim about which component contributes more to classification
performance.

\textbf{Source.} Direct arithmetic over the implementation's module counts.
No paper is claimed.

\subsubsection{Sampler weight}
\label{fg:sampler-ch4}
\textbf{Thesis location:} 4.7.2, \texttt{sec:sampler-config}.

\textbf{No new formula.} Chapter~4 states the \texttt{WeightedRandomSampler}
weight $1/n_c$, which is Equation~\eqref{fg:eq:sampling}.

Chapter~4 fixes the scope: $n_c$ counts training examples of that example's
class \emph{within the current fold's training split}, sampling proceeds with
\texttt{replacement=True}, and each epoch draws as many indices as the split
contains. This is also what makes the class-weighting rule well founded rather
than merely convenient --- because the skew correction genuinely happens at
the sampler, disabling its loss-side equivalent avoids double-counting rather
than removing the correction.

\textbf{Source.} As Equation~\eqref{fg:eq:sampling}: project implementation.

\subsection{Metric selection: thesis section 4.9}

\subsubsection{The mean-of-folds macro F1 ceiling}
\label{fg:fold-ceiling}
\textbf{Thesis location:} 4.9.3, \texttt{sec:mean-of-folds-rejected}.
% SOURCE-DISPLAY-CH4-02
\begin{equation}
	\frac{10 \cdot \tfrac13 + 8 \cdot \tfrac23 + 7 \cdot 1}{25} = \frac{3.333\ldots + 5.333\ldots + 7}{25} \approx 0.627.
    \label{fg:eq:fold-ceiling}
\end{equation}
The three numerator terms group the 25 leave-one-subject-out folds by how many
of the three classes appear in each fold's held-out ground truth:
\textbf{10} folds contain one class, \textbf{8} contain two, \textbf{7}
contain all three. The fractions $\tfrac13$, $\tfrac23$ and $1$ are the
maximum macro F1 each group can attain.

The bound follows from the zero convention recorded at
Equation~\eqref{fg:eq:metrics}. A macro F1 computed \emph{within} a fold is
the unweighted mean of that fold's three per-class F1 scores, and a class
absent from that fold's ground truth scores zero no matter what the model
predicts. A single-class fold is therefore capped at $1/3$ and a two-class
fold at $2/3$. Averaging across folds averages seventeen sub-unity ceilings
into every score. These ceilings are fixed by which subject was held out,
before any model is trained, and apply identically to all twelve
configurations.

Pooling avoids the ceiling entirely, which is the argument of
Equation~\eqref{fg:eq:pooling}: accumulating $TP$, $FP$ and $FN$ across all 25
folds and all 156 clips before computing any per-class F1 means a class
locally absent from one fold no longer forces a zero into its contribution.
The consequence is arithmetic --- a metric with a hard ceiling below 0.7
cannot serve as a headline result --- so pooled macro F1 is primary and the
stored \texttt{macro\_f1} field is not. No run artefact stores the pooled
quantity under any key; it is derived afterwards by the reporting scripts
under \texttt{tools/}.

The same fold composition biases mean-of-folds \emph{accuracy} in the same
direction without imposing a hard ceiling, since one predicted label suffices
to score well on a single-class fold. For the full configuration, mean-of-folds
accuracy exceeds pooled accuracy by \textbf{6.3 points}.

\textbf{Source.} This thesis's own arithmetic over its fold composition. It is
not a published result. See et al.~\cite{see2019} supply the pooled
construction it argues for.


\section{Chapter 5: Results}
\label{fg5:chapter5}

\textbf{Source file:} \texttt{contents/chapter5\_shortened.tex}. The chapter
contains \textbf{no displayed equation blocks and defines no new formula}. What
it contains is fourteen inline mathematical spans, and all but two of those are
notation rather than formulas: tensor shapes, configuration-pair arrows and a
difference symbol.

That is what a results chapter should look like. Its work is to report measured
numbers, and the formulas that produce them were defined in Chapters 2 to 4.
This section therefore does three things, none of which restates an earlier
formula: it promotes the two aggregation rules the chapter uses throughout but
only ever states in prose; it works through the one genuine calculation the
chapter performs, the all-Negative reference classifier, as a numeric
evaluation of metric formulas defined earlier; and it indexes the chapter's
notation.

\subsection{Aggregation rules used throughout the chapter}

\subsubsection{The matched-pair difference}
\label{fg5:pair-difference}
\textbf{Thesis locations:} 5.2--5.5, every ``The six pairs'' and ``The four
pairs'' table.

The chapter's central quantity is written in its tables as
``$\Delta$ pooled macro F1'' and in its prose as a signed change, but never as
an equation. Promoted, it is
\begin{equation}
    \Delta_j = S\!\left(\mathrm{on}_j\right) - S\!\left(\mathrm{off}_j\right),
    \label{fg5:eq:pair-difference}
\end{equation}
where $j$ indexes a matched pair, $S(\cdot)$ is pooled macro F1 as constructed
in \eqref{fg:eq:macro} and pooled per \eqref{fg:eq:pooling}, and
$\mathrm{on}_j$ and $\mathrm{off}_j$ are the two configurations of pair $j$ ---
identical in every flag except the one under test.

The construction is what gives the chapter its evidential standing. Because
the two members differ in exactly one flag, and because every non-varied
factor, seed and fold assignment is held identical, $\Delta_j$ isolates that
flag within that architectural context. The chapter's arrow notation,
$\texttt{config\_1} \rightarrow \texttt{config\_2}$, names the pair by its off
and on members in that order.

Two cautions the chapter states explicitly. Differences are calculated from
\emph{unrounded} results, so subtracting the displayed four-decimal scores can
disagree in the last digit. And a single seed means $\Delta_j$ is a point
estimate carrying no uncertainty interval; it cannot establish statistical
significance.

\textbf{Source.} Project design. The matched-pair form is the study's own
construction, not an aggregation rule imported from a reviewed paper.

\subsubsection{The mean marginal effect}
\label{fg5:mean-effect}
\textbf{Thesis locations:} 5.2.1, 5.3.1, 5.4.1, 5.5.1, each ``The effect''.

Each component's headline number is the mean of its pair differences:
\begin{equation}
    \overline{\Delta} = \frac1J \sum_{j=1}^{J} \Delta_j,
    \label{fg5:eq:mean-effect}
\end{equation}
with $J = 6$ for the transformer and for EVM, and $J = 4$ for SimAM and for
the convolutional stem. The counts differ because the validity rule of
\eqref{fg:eq:factorial} removes the four cells that pair SimAM with no CNN, so those
two components have fewer architecturally valid pairs to toggle within.

The chapter is careful about what this average can carry, and the guide
repeats that care rather than the algebra. A mean over four or six
single-seed differences has no distributional interpretation: it does not
estimate a population effect, and its spread across pairs is evidence about
architectural dependence rather than sampling noise. Chapter~5's own
subsection ``Why the mean is not the finding'' makes this argument for SimAM,
where the sign varies across pairs.

For the transformer, $\overline{\Delta} = +0.2173$ with all six differences
positive, ranging $+0.1578$ to $+0.2786$. Consistency of sign across six of
six pairs is a stronger observation than the mean itself, because it does not
depend on the magnitudes.

\textbf{Source.} Project design, as above.

\subsection{The all-Negative reference classifier: thesis section 5.1.3}

\subsubsection{Numeric evaluation of the metric definitions}
\label{fg5:all-negative}
\textbf{Thesis location:} 5.1.3, \texttt{sec:all-negative-reference}.

\textbf{No new formula.} This is the chapter's one genuine calculation, and it
evaluates formulas already defined: accuracy via
\eqref{fg:eq:micro}, per-class F1 via \eqref{fg:eq:metrics} and the macro average via
\eqref{fg:eq:macro}. Only the arithmetic is new.

Consider a classifier that ignores its input and predicts Negative for every
one of the 156 clips, 99 of which are Negative. Then
\begin{equation}
\begin{aligned}
    \text{accuracy} &= \frac{99}{156} = 0.6346,\\
    F1_{\mathrm{Negative}} &= \frac{2\times 99}{156+99} = 0.7765,\\
    F1_{\mathrm{macro}} &= \frac{0.7765+0+0}{3} = 0.2588.
    \end{aligned}
    \label{fg5:eq:all-negative}
\end{equation}

The middle line is \eqref{fg:eq:metrics} with $TP = 99$, $FP = 57$ and $FN = 0$:
every Negative clip is found, so there are no false negatives, and every
non-Negative clip is wrongly called Negative, giving
$2TP + FP + FN = 198 + 57 = 255 = 156 + 99$. The third line is the macro
average with the Positive and Surprise scores at zero, since the classifier
never predicts either, and the project convention of \eqref{fg:eq:metrics} assigns
zero rather than omitting an undefined class.

\textbf{What it establishes, and what it does not.} All twelve configurations
exceed 0.2588 on pooled macro F1. That is a real but weak reference: it
rules out one specific degenerate rule, not every input-independent rule, and
it does not establish generalisation. The accuracy line carries the sharper
warning --- \textbf{0.6346 accuracy is attainable without discriminating among
expressions at all}, which is why accuracy is not the primary metric and why
configuration~8's accuracy lead over configuration~2 (0.7500 against 0.7436,
117 against 116 correct clips) does not settle the ranking.

\textbf{Source.} This thesis's own arithmetic over its class distribution. The
metric definitions are attributed at their own equations; See et
al.~\cite{see2019} supply the pooled construction.

\subsubsection{The mean-of-folds ceiling, restated}
\label{fg5:ceiling-ch5}
\textbf{Thesis location:} 5.1.3.

\textbf{No new formula.} Chapter~5 reports a \textbf{0.6267} ceiling on
mean-of-folds macro F1 imposed by fold composition. This is
\eqref{fg:eq:fold-ceiling}, which Chapter~4 states as $\approx 0.627$; the two agree, the
results chapter simply carrying one more digit. The ceiling applies with all
three classes retained and undefined class F1 set to zero, and it does not
apply to pooled macro F1. Figure 5.2 of the thesis plots the two aggregations
against each other with that ceiling marked.

\textbf{Source.} As \eqref{fg:eq:fold-ceiling}: this thesis's own arithmetic over its fold
composition.

\subsection{Notation index}
\label{fg5:notation}

The chapter's remaining mathematical spans are notation, not formulas. They
are listed here so the index is complete.

\textbf{$\rightarrow$ (configuration pairs).} In
$\texttt{config\_1} \rightarrow \texttt{config\_2}$ and the abbreviated forms
$2\rightarrow12$, $9\rightarrow7$, $5\rightarrow16$, $13\rightarrow16$,
$9\rightarrow6$ and $2\rightarrow9$, the arrow separates the off member from
the on member of a matched pair, in the sense of
Equation~\eqref{fg5:eq:pair-difference}. The numerals are configuration
identifiers, not ranks or magnitudes.

\textbf{$\Delta$.} The difference of
Equation~\eqref{fg5:eq:pair-difference}, used as a column heading.

\textbf{$-$.} A typeset minus sign on negative differences, written
explicitly so that a leading minus is not mistaken for a hyphen.

\textbf{$(1,3,3)$.} The convolutional kernel of the stem, in the
$(k_T,k_H,k_W)$ convention set out in Chapter 2's entry on
three-dimensional tensor and kernel notation. Because $k_T = 1$, the stem performs \textbf{no temporal
convolution}; Chapter~5 relies on this when it describes the transformer-free
control as order-blind. The control is nevertheless not independent across
frames at every operation, since batch normalisation during training and
SimAM whenever enabled both use statistics spanning time.

\textbf{$224\times224$ and $28\times28$.} Spatial input dimensions in pixels,
in the width-by-height pixel convention of Chapter 2's image-dimensions
entry. The comparison is raised as one of three unresolved candidate
explanations for the stem's weak contribution: this pipeline's
$224\times224$ input is large relative to STSTNet's
$28\times28$~\cite{liong2019b}. The chapter declines to conclude from it,
noting that other reviewed systems also use $224\times224$ and that shallow
models need not share deeper models' resolution sensitivity. No universal
literature threshold is established.

\subsection{Coverage boundary}

All fourteen inline mathematical spans in the shortened Chapter 5 are
accounted for above: two promoted aggregation rules, one numeric evaluation
across three lines, one restated ceiling, and ten notation items. The chapter
has \textbf{no displayed equation blocks}, so none has been omitted, and no
formula from Chapters 1 to 4 has been reprinted here.

Quantities the chapter reports without expressing as a calculation --- pooled
macro F1 scores, peak VRAM, extrapolated GPU-hours, per-class F1 tables and
confusion matrices --- are measurements, not formulas, and have not been
reverse-engineered into equations.

\section{Chapter 6: Conclusion}
\label{fg6:chapter6}

\textbf{Source file:} \texttt{contents/chapter6.tex}. No shortened version of
that chapter exists. It contains \textbf{no displayed mathematics, no new
formula, and no restatement of an earlier formula in algebraic form}. It also
contains no citations at all.

What it contains is nine inline mathematical spans, every one of which is a
\emph{reference to a quantity already fixed elsewhere} --- a design-space
count, four inherited hyper-parameters, one weighting ratio and a set of image
dimensions. A conclusion that introduced new mathematics would be doing
something wrong; this one does not. This section therefore indexes those nine
spans: what each denotes, what value it carries, which equation above defines
it, and --- the part that matters for reading the conclusion --- what claim the
chapter attaches to it.

\subsection{The design-space count}

\subsubsection{The \texorpdfstring{$2^4$}{2\^{}4} matrix}
\label{fg6:factorial}
\textbf{Thesis location:} 6.1, Contributions, first bullet.

\textbf{No new formula.} The chapter writes ``every architecturally valid cell
of a $2^4$ matrix''. This is \eqref{fg:eq:factorial}, which gives $2^4 = 16$ cells of
which 12 are architecturally valid, the four invalid ones being those pairing
SimAM with no convolutional stem.

The conclusion uses the count to state the study's form: twelve configurations,
each trained from scratch, each scored under complete 25-fold
leave-one-subject-out, with every non-varied factor held identical. The
qualifier ``architecturally valid'' is doing real work --- it is what
distinguishes twelve from sixteen, and \eqref{fg:eq:factorial} records that the four
excluded cells are rejected as ill-formed specifications rather than as
weaker pipelines.

\subsection{The inherited preprocessing parameters}
\label{fg6:inherited}
\textbf{Thesis location:} 6.2, Limitations, first paragraph; and 6.3, Further
Work, the divergence backlog.

\textbf{No new formulas.} The chapter names four fixed parameters, all
defined earlier and none introduced here. They are gathered in one place
because the conclusion's central caveat is about them jointly: they are
\emph{inherited from the literature, held fixed throughout, and never swept},
even though the reviewed work establishes that such parameters have interior
optima locatable only by exhaustive search.

\begin{center}
\begin{tabular}{@{}l l l p{5.4cm}@{}}
\hline
\textbf{Symbol} & \textbf{Value} & \textbf{Defined at} & \textbf{What it controls} \\
\hline
$\alpha$  & 10        & \eqref{fg:eq:evm-gain}  & Eulerian magnification gain: the factor by which band-passed temporal variation is amplified before reconstruction \\
$T$       & 32        & Ch.\ 2, tensor shape & Temporal length of the motion tensor --- 32 flow fields from 33 sampled frames \\
$\lambda$ & $10^{-4}$ & \eqref{fg:eq:simamgate} & SimAM regularisation constant, stabilising the division by channel variance \\
$\gamma$  & 2.0       & \eqref{fg:eq:focal} & Focal-loss focusing exponent, controlling how strongly confident examples are down-weighted \\
\hline
\end{tabular}
\end{center}

\textbf{The claim attached to them.} Every effect reported in Chapter~5 is
\emph{conditional on these four values}. The design measures architecture at
one point in preprocessing space and cannot say how any effect would move if
that point changed. The conclusion lists sweeping them --- $T$ on one hand,
and $\alpha$, $\lambda$ and $\gamma$ on the other --- among the further work
that would establish whether the reported effects survive away from the single
setting at which they were measured.

\textbf{A notational warning carried over.} These symbols are reused across
topics with unrelated meanings. $\alpha$ is magnification gain here but the
class-weight factor $\alpha_t$ in \eqref{fg:eq:focal}; $\gamma$ is the focal
exponent here but the per-channel scale $\gamma_c$ in \eqref{fg:eq:bn};
$\lambda$ appears only in SimAM. They are not interchangeable, and the
conclusion's single sentence naming all four together is the one place in the
thesis where confusing them is easiest.

\subsection{The strain shear-weighting ratio}

\subsubsection{\texorpdfstring{$\sqrt{2}$}{sqrt(2)}}
\label{fg6:sqrt2}
\textbf{Thesis location:} 6.3, Further Work, the declared-divergence backlog.

\textbf{No new formula.} The chapter refers to ``the four-term strain
magnitude, whose $\sqrt{2}$ shear weighting the three-term form \ldots\ does
not carry''. This ratio is the consequence of Equations~\eqref{fg:eq:frobenius}
and~\eqref{fg:eq:strain-magnitude}: the published four-term magnitude counts shear twice and the
implemented three-term magnitude counts it once, so for pure shear the two
scalars stand in the ratio $\sqrt{2}$.

The conclusion's use of it is to schedule an experiment, not to restate the
algebra: adopting the four-term magnitude is listed in the backlog of declared
divergences worth testing. Chapter~3 establishes that the two conventions were
never compared experimentally, so the ratio marks an untested difference
rather than a known cost.

\subsection{Image dimensions}
\label{fg6:dimensions}
\textbf{Thesis locations:} 6.2, Limitations, second paragraph; 6.3, Further
Work, second paragraph.

\textbf{No new formula.} The chapter names four spatial resolutions, in the
width-by-height pixel convention defined in Chapter~2's geometry entry:
the implemented $224 \times 224$; the literature norm it contrasts that with,
``at or below $100 \times 100$''; and the two resolutions it proposes
re-training the stem arm at, $28 \times 28$ and $56 \times 56$.

\textbf{The claim attached to them.} This is the conclusion's sharpest
qualification of a Chapter~5 result. Because the convolutional stem runs far
above the literature's resolution range, its measured negative effect
\emph{may be a resolution result as much as an architectural one}. The chapter
advances that reading and then qualifies it, noting that the same source
reports shallow models to be largely robust to resolution. Re-training the
stem arm at $28 \times 28$ and $56 \times 56$ with everything else held as
Chapter~4 specifies is named as the cheapest experiment left undone, and the
one whose outcome would most change how the stem result must be worded.

\subsection{Coverage boundary}

All nine inline mathematical spans in Chapter 6 are accounted for above. The
chapter has \textbf{no displayed equation blocks} and \textbf{no citations},
so none of either has been omitted.

Three further quantities appear in the chapter as \emph{reported measurements}
rather than as formulas, and have not been reverse-engineered into equations:
the two headline pooled macro F1 scores, \textbf{0.7122} for
\texttt{config\_2\_temporal\_only} against \textbf{0.6659} for
\texttt{config\_8\_proposed\_unified}; the count of twelve declared
divergences; and the compute share attributed to the stem-bearing
configurations. Each is produced by a formula defined in Chapters 2 to 4 ---
pooled macro F1 at \eqref{fg:eq:macro}, pooled per \eqref{fg:eq:pooling} --- and each is
a number this chapter quotes rather than derives.

The conclusion also proposes one statistical procedure it does not formalise:
McNemar's test over the clips on which two configurations disagree, offered as
what retained per-clip predictions would make possible. No equation for it
appears in the thesis, and none has been invented here.

\section{Coverage boundary for this guide}
\label{fg:coverage}

This is the complete inventory. The thesis contains \textbf{25 displayed
equation blocks}; this guide prints \textbf{40 numbered equations}, the
difference being the inline quantities and prose relations it promotes, less
the restatements it cross-references rather than repeats.

\textbf{Chapter 1} contains no formulas; this is stated rather than remedied by
importing equations from later chapters.

\textbf{Chapter 2} contributes all sixteen of its displayed blocks, plus its
inline formulas promoted as explanatory expansions.

\textbf{Chapter 3} contains seven displayed blocks. Two are printed here ---
the four-term published magnitude at Equation~\eqref{fg:eq:frobenius} and the
shear identity at Equation~\eqref{fg:eq:shear-identity}. The remaining five
restate Chapter 2 relations and are cross-referenced instead: the per-class F1
and its unweighted average (\autoref{fg:uf1}), the implemented strain
magnitude and its shear component (\autoref{fg:strain-divergence}), the SimAM
statistics and gate (\autoref{fg:simam-ch3}), and focal loss
(\autoref{fg:focal-ch3}). Three of its inline quantities are promoted:
Equations~\eqref{fg:eq:flow-magnitude}, \eqref{fg:eq:effective-rate}
and~\eqref{fg:eq:evm-interval}.

\textbf{Chapter 4} contains two displayed blocks, both printed here:
Equations~\eqref{fg:eq:posenc} and~\eqref{fg:eq:fold-ceiling}. Two of its
inline quantities are promoted, Equations~\eqref{fg:eq:factorial}
and~\eqref{fg:eq:param-totals}; its sampler weight restates
Equation~\eqref{fg:eq:sampling} and is cross-referenced at
\autoref{fg:sampler-ch4}.

\textbf{Chapter 5} contains no displayed blocks. Its fourteen inline spans are
all accounted for: two promoted aggregation rules
(Equations~\eqref{fg5:eq:pair-difference} and~\eqref{fg5:eq:mean-effect}), one
numeric evaluation of earlier metric definitions
(Equation~\eqref{fg5:eq:all-negative}), one restated ceiling
(\autoref{fg5:ceiling-ch5}) and ten notation items
(\autoref{fg5:notation}).

\textbf{Chapter 6} contains no displayed blocks and no citations. Its nine
inline spans are all back-references to quantities fixed in Chapters 2 to 4,
indexed at \autoref{fg6:chapter6}. It introduces no formula of its own.

\textbf{What has deliberately not been done.} No formula has been invented to
fill a gap, no equation has been reverse-engineered from a reported
measurement, and no originating paper has been cited merely because a technique
shares a name with one. Quantities the thesis reports as measurements ---
scores, parameter counts, VRAM, GPU-hours, confusion matrices --- are
measurements, not formulas, and appear here only where the thesis itself shows
their arithmetic.


\begin{thebibliography}{99}
\bibitem{yan2014}
W.-J. Yan et al., ``CASME II: An Improved Spontaneous Micro-Expression Database
and the Baseline Evaluation,'' \emph{PLoS ONE}, vol. 9, no. 1, e86041, 2014.

\bibitem{bai2021}
Mengjiong Bai, Roland Goecke and Damith Herath, ``Micro-Expression Recognition
Based on Video Motion Magnification and Pre-Trained Neural Network,''
\emph{IEEE International Conference on Image Processing}, pp. 549--553, 2021.

\bibitem{liong2019a}
Y. S. Gan, Sze-Teng Liong, Wei-Chuen Yau, Yen-Chang Huang and Lit Ken Tan,
``OFF-ApexNet on Micro-Expression Recognition System,''
\emph{Signal Processing: Image Communication}, vol. 74, pp. 129--139, 2019.
The supplied preprint lists Liong first; equation locators above refer to that PDF.

\bibitem{zhaoS2021}
Sirui Zhao et al., ``A Two-Stage 3D CNN Based Learning Method for Spontaneous
Micro-Expression Recognition,'' \emph{Neurocomputing}, vol. 448,
pp. 276--289, 2021.

\bibitem{shreve2011}
Matthew Shreve, Sridhar Godavarthy, Dmitry Goldgof and Sudeep Sarkar,
``Macro- and Micro-Expression Spotting in Long Videos Using Spatio-Temporal
Strain,'' \emph{IEEE International Conference on Automatic Face and Gesture
Recognition}, 2011.

\bibitem{yang2021}
Lingxiao Yang, Ru-Yuan Zhang, Lida Li and Xiaohua Xie,
``SimAM: A Simple, Parameter-Free Attention Module for Convolutional Neural
Networks,'' \emph{International Conference on Machine Learning}, 2021.

\bibitem{dosovitskiy2021}
Alexey Dosovitskiy et al., ``An Image Is Worth $16\times16$ Words:
Transformers for Image Recognition at Scale,''
\emph{International Conference on Learning Representations}, 2021.

\bibitem{zhang2022}
Liangfei Zhang, Xiaopeng Hong, Ognjen Arandjelovi{\'c} and Guoying Zhao,
``Short and Long Range Relation Based Spatio-Temporal Transformer for
Micro-Expression Recognition,'' \emph{IEEE Transactions on Affective
Computing}, vol. 13, no. 4, pp. 1973--1985, 2022.

\bibitem{see2019}
John See, Moi Hoon Yap, Jingting Li, Xiaopeng Hong and Su-Jing Wang,
``MEGC 2019---The Second Facial Micro-Expressions Grand Challenge,''
\emph{IEEE International Conference on Automatic Face and Gesture
Recognition}, 2019.
\bibitem{liong2014a}
Sze-Teng Liong, Raphael C.-W. Phan, John See, Yee-Hui Oh and KokSheik Wong,
``Optical Strain Based Recognition of Subtle Emotions,'' \emph{International
Symposium on Intelligent Signal Processing and Communication Systems
(ISPACS)}, 2014.

\bibitem{liong2014b}
Sze-Teng Liong, John See, Raphael C.-W. Phan, Anh Cat Le Ngo, Yee-Hui Oh and
KokSheik Wong, ``Subtle Expression Recognition Using Optical Strain Weighted
Features,'' \emph{Computer Vision -- ACCV 2014 Workshops}, Lecture Notes in
Computer Science, vol. 9009, pp. 644--657, Springer, 2015.

\bibitem{liong2016}
Sze-Teng Liong, John See, Raphael C.-W. Phan, Yee-Hui Oh, Anh Cat Le Ngo,
KokSheik Wong and Su-Wei Tan, ``Spontaneous Subtle Expression Detection and
Recognition Based on Facial Strain,'' \emph{Signal Processing: Image
Communication}, vol. 47, pp. 170--182, 2016.

\bibitem{liong2019b}
Sze-Teng Liong, Y. S. Gan, John See, Huai-Qian Khor and Yen-Chang Huang,
``Shallow Triple Stream Three-Dimensional CNN (STSTNet) for Micro-Expression
Recognition,'' \emph{14th IEEE International Conference on Automatic Face and
Gesture Recognition (FG 2019)}, 2019.

\end{thebibliography}
\end{document}
